% F15 -- Levelling every user to 200 records: what the score looks like, and what it is.
\begin{figure}[t]
\centering
\replegend{\swatch{clbase}~persistence baseline \qquad \swatch{clmodel}~model}
\fitwidth{%
\begin{tikzpicture}
\begin{groupplot}[
  group style={group size=2 by 1, horizontal sep=2.2cm, ylabels at=edge left},
  repbar, width=0.42\linewidth, height=5.4cm, xtick=data,
  title style={font=\footnotesize,yshift=1pt},
]
% ---- (a) each run against its own baseline ----
\nextgroupplot[title={(a) each run against \emph{its own} baseline},
  symbolic x coords={nat,lev},
  xticklabels={natural lengths,levelled to 200},
  ymin=0, ymax=122, ytick={0,20,...,80},
  ylabel={top-1 accuracy (\%)}, /pgf/bar width=16pt,
  x tick label style={font=\scriptsize}]
\addplot[fill=clbase,draw=none]  coordinates {(nat,43.1)(lev,76.5)};
\addplot[fill=clmodel,draw=none] coordinates {(nat,46.8)(lev,74.9)};
\node[font=\small,text=clgain,anchor=south] at (axis cs:nat,60) {$+3.7\pt$};
\node[font=\small,text=clloss,anchor=south] at (axis cs:lev,102) {$-1.6\pt$};
%
% ---- (b) the arithmetic of copying ----
\nextgroupplot[title={(b) what the 4{,}000 scored positions are},
  symbolic x coords={cop,real},
  xticklabels={{copies\\2{,}262 of 4{,}000},{real decisions\\1{,}738 of 4{,}000}},
  ymin=0, ymax=125, ytick={0,25,...,100},
  /pgf/bar width=16pt,
  x tick label style={font=\scriptsize,align=center}]
\addplot[fill=clbase,draw=none]  coordinates {(cop,100.0)(real,45.9)};
\addplot[fill=clmodel,draw=none] coordinates {(real,42.2)};
\node[font=\scriptsize,text=clloss,anchor=south,xshift=-11pt] at (axis cs:real,58) {$-3.7\pt$};
\end{groupplot}
\end{tikzpicture}}
\caption{Levelling every user to 200 records}
\label{fig:resampling}
\figtext{%
Stretching every sequence to exactly 200 records makes the score look
spectacular and turns the gap negative. (a)~Each run against \emph{its own}
persistence baseline, natural lengths on the left and levelled sequences on
the right; only the comparison inside a pair means anything. The levelled
model reaches 74.9\,\%, twenty-eight points above the natural-length run and
by a wide margin the highest absolute accuracy in the whole project, but the
baseline on those same sequences reaches 76.5\,\% and the gap changes sign,
from $+3.7\pt$ to $-1.6\pt$. (b)~The reason is arithmetic rather than
modelling. Of the 4{,}000 scored positions (100 users $\times$ the last 40 of
their 200 records), 2{,}262 ask for a record copied from the one before, and
any rule that repeats is right on all of them by construction; both the model
and the baseline are handed those free answers, which is why both rose by
about thirty points. On the 1{,}738 positions that are real decisions the
baseline scores 45.9\,\% against the model's 42.2\,\%. The technique does
deliver what it promises, the longest-to-shortest ratio falling from 200:19 to
1:1 and the Gini coefficient of per-user weight to zero, but a copied step
never changes site, so levelling compresses the between-user spread in site-change rate by a factor of $2.3$ and drags every user towards the homebody
profile, which disables the group detection the project's only real gain
depends on. The model bar in panel~(b) is derived from the reported 74.9\,\%,
not re-measured: it assumes the model never misses a copy, and rises to
44.8\,\% if it misses 2\,\% of them, both ends staying below the baseline.}
\end{figure}
